% =============================================================================
% CATS Group 9 - Draft Research Paper
% Breast cancer subtype classification from aCGH copy-number profiles.
% Bioinformatics journal style (OUP).
% =============================================================================
%
% To compile:  make        (or: pdflatex main && bibtex main && pdflatex main && pdflatex main)
% To watch:    make watch  (requires latexmk)
%
% =============================================================================

\documentclass[11pt,a4paper,twocolumn]{article}

% ---- Packages ----
\usepackage[utf8]{inputenc}
\usepackage[T1]{fontenc}
\usepackage{lmodern}
\usepackage{microtype}
\usepackage[margin=2cm]{geometry}
\usepackage{graphicx}
\usepackage{booktabs}
\usepackage{amsmath}
\usepackage{amssymb}
\usepackage{xcolor}
\usepackage{hyperref}
\usepackage[round,authoryear]{natbib}
\usepackage{caption}
\usepackage{subcaption}
\usepackage{float}
\usepackage{enumitem}
\usepackage{multicol}
\usepackage{authblk}

% ---- Figure path ----
\graphicspath{{figures/}}

% ---- Hyperref setup ----
\hypersetup{
  colorlinks=true,
  linkcolor=blue!60!black,
  citecolor=blue!60!black,
  urlcolor=blue!60!black
}

% ---- Bibliography style ----
\bibliographystyle{abbrvnat}

% ---- Title ----
\title{Hierarchical classification of breast cancer molecular subtypes
       from discrete aCGH copy-number profiles}

% ---- Authors ----
\author[1]{Alexandros Michailidis}
\author[1]{Antonie Wagner}
\author[1]{Christos Botos}
\author[1]{Yan Qiao}
\affil[1]{Computer Science and Bioinformatics Master's Programmes,
          Vrije Universiteit Amsterdam, The Netherlands}

\date{}

% =============================================================================
% WRITING GUIDE (read before expanding any section into prose):
%
%   - Mixed active/passive voice. Past tense for what was done, present for
%     what is generally true.
%   - Every paragraph: CLAIM sentence first, then EVIDENCE/DETAIL.
%   - Precise numbers, no "several", "many", "good".
%   - NO dashes as parenthetical asides. Use commas or parentheses.
%   - Citations integrated into argument, not dropped at end of sentence.
%   - Methods must be reproducible: grids, seeds, software versions.
%   - Acknowledge uncertainty with reasons, not hedge words.
% =============================================================================

\begin{document}
\maketitle

% =============================================================================
\begin{abstract}
\noindent
\textbf{Motivation:}
Breast cancer subtypes (HER2+, HR+, TN) require different treatments.
Array CGH provides discrete copy-number profiles for molecular classification.
Wessels et al.\ (2005) showed simple classifiers match complex ones on expression data;
unknown whether this holds for discrete aCGH with limited samples ($n=100$).

\textbf{Results:}
Flat 2x2 experiment: RF outperforms NMC, apparently contradicting Wessels.
HER2+ perfectly separable (masks real HR+/TN problem).
Hierarchical decomposition: NMC outperforms RF for HR+/TN (reversal).
Wessels' principle is conditional on multi-class structure.
Best biomarker: chr17q12 ERBB2 amplicon.
Final model: hierarchical EN+NMC with plateau ensembling, Stage~2 BA $= 0.76$, combined 3-class BA $= 0.84$.
Validation estimate: 48/57 correct (per-class recall weighted by training proportions).

\textbf{Availability:}
Code and predictions in submission archive.
\end{abstract}

% =============================================================================
\section{Introduction}
\label{sec:introduction}

\begin{itemize}[leftmargin=*, nosep]
  \item Breast cancer is heterogeneous; molecular subtypes drive treatment
  \item HER2+, HR+, Triple Negative -- different prognosis and therapy
  \item \cite{perou2000molecular} intrinsic subtypes; \cite{sorlie2001gene} clinical implications; \cite{slamon1987her2} HER2 amplification
  \item Array CGH measures genome-wide CN changes \citep{pinkel2005array}
  \item Discrete CN calls: loss ($-1$), normal ($0$), gain ($+1$), amplification ($+2$)
  \item Preprocessing: CBS segmentation \citep{venkatraman2007cbs}, CGHcall \citep{vandewiel2007cghcall}, overview \citep{vandewiel2011preprocessing}
  \item Goal: classify subtypes from CN profiles alone (n=100 train, 57 validation)
  \item \cite{wessels2005protocol}: simple classifiers + univariate filtering match complex on expression data
  \item Tested on van 't Veer (2002) breast cancer data; key insight: complex models overfit at small n
  \item \cite{haury2011influence}: univariate filters competitive, stability matters
  \item Gap: Wessels used continuous expression, not discrete aCGH
  \item Gap: multi-class structure effect on simple-vs-complex balance untested
  \item \textbf{RQ:} Does Wessels' principle hold for discrete aCGH classification, and how do feature redundancy and multi-class structure affect classifier selection?
  \item Approach: 2x2 factorial design, nested CV, hierarchical extension
  \item Preview: the answer depends on problem decomposition
\end{itemize}

% =============================================================================
\section{Methods}
\label{sec:methods}

\subsection{Data description}
\label{sec:data}

\begin{itemize}[leftmargin=*, nosep]
  \item 100 labelled samples, 244K-probe aCGH platform
  \item Preprocessed into 2834 genomic regions with discrete CN calls
  \item Classes: HER2+ ($n=32$), HR+ ($n=36$), TN ($n=32$) -- approximately balanced
  \item 57 unlabelled validation samples, same platform
  \item Coordinates: hg18/NCBI Build 36
  \item Implemented in Python 3.10, scikit-learn \citep{pedregosa2011sklearn}
\end{itemize}

\subsection{Label-free region merging}
\label{sec:merging}

\begin{itemize}[leftmargin=*, nosep]
  \item Adjacent regions highly correlated: median $r = 0.92$ within 0.1 Mb
  \item Redundancy inflates feature space without adding discriminative info
  \item Inspired by CGHregions \citep{vanwieringen2007cghregions} -- same type of correlation-based reduction for aCGH
  \item Algorithm: greedy one-pass per chromosome, merge adjacent if Pearson $r > 0.8$
  \item Representation: median CN of constituent regions per sample
  \item Threshold justification: 91\% of adjacent pairs exceed $r = 0.8$; only 4.4\% below $r = 0.7$
  \item Result: 2834 regions $\to$ 273 segments (90.4\% reduction)
  \item Signal preservation: top KW features, PCA structure, t-SNE topology unchanged
  \item Sensitivity: BA varies $<2$ pp across $r \in \{0.7, 0.8, 0.9\}$; no merging $-7$ pp worse ($p = 0.008$, Wilcoxon signed-rank, directional only)
  \item \textbf{Leakage defense:} label-free (no class info used), applied before CV
  \item Correlation structure of 80 samples $\approx$ 100 samples given strong autocorrelation
  \item Consistent with CGHregions global preprocessing; distinct from \cite{ambroise2002selection} label-dependent leakage
  \item Merging map saved in model artifact, applied to validation data at prediction time
\end{itemize}

\subsection{Experimental design}
\label{sec:design}

\begin{itemize}[leftmargin=*, nosep]
  \item $2 \times 2$ factorial: (KW, EN) feature selection $\times$ (NMC, RF) classifier
  \item Directly tests \cite{wessels2005protocol} recommendation
\end{itemize}

\subsubsection{Feature selection}
\label{sec:features}

\begin{itemize}[leftmargin=*, nosep]
  \item \textbf{KW:} Kruskal-Wallis $H$-statistic ranking \citep{kruskalwallis1952}, select top $k$
  \item Non-parametric, respects ordinal CN values, handles 3-class natively
  \item $k \in \{5, 10, 15, 20, 30, 50, 75, 100\}$ tuned via inner CV
  \item \textbf{EN:} Elastic net multinomial logistic regression \citep{zou2005elasticnet}
  \item L1+L2 penalty: sparsity + grouping effect for correlated features
  \item Select top $k \in \{5, 10, 20, 50\}$ by absolute coefficient
  \item $C \in \{10^{-3}, \ldots, 10^{2}\}$ (10 log-spaced), $\lambda_1 \in \{0.1, 0.3, 0.5, 0.7, 0.9\}$
  \item Inner grid: $10 \times 5 \times 4 = 200$ configs per fold
\end{itemize}

\subsubsection{Classifiers}
\label{sec:classifiers}

\begin{itemize}[leftmargin=*, nosep]
  \item \textbf{NMC:} Nearest mean classifier -- assign to closest class centroid (Euclidean)
  \item Linear boundary, zero base hyperparams; shrinkage $\delta \in \{\text{None}, 0.1, 0.2, 0.5\}$ tuned
  \item Referenced by Wessels as prototypical simple classifier
  \item \textbf{RF:} Random forest, 500 trees \citep{breiman2001random}, balanced class weights
  \item Tuned: max\_features $\in \{\text{sqrt}, 0.3\}$, min\_samples\_leaf $\in \{1, 5\}$ (4 configs)
  \item Non-linear, captures feature interactions
\end{itemize}

\subsubsection{The four pipelines}
\label{sec:pipelines}

\begin{itemize}[leftmargin=*, nosep]
  \item KW+NMC: $8 \times 4 = 32$ inner configs
  \item KW+RF: $8 \times 4 = 32$ inner configs
  \item EN+NMC: $200 \times 4 = 800$ inner configs
  \item EN+RF: $200 \times 4 = 800$ inner configs
  \item Grid size asymmetry relevant for plateau ensembling interpretation
\end{itemize}

\subsection{Hierarchical classifier}
\label{sec:hierarchical}

\begin{itemize}[leftmargin=*, nosep]
  \item Motivation: flat experiment $\to$ HER2+ recall = 1.00 in all pipelines; bottleneck is HR+/TN confusion
  \item Dedicated HR+/TN KW ranking surfaces different features: chr12q (23/30), chr5q (7/30) instead of chr17 ERBB2
  \item \textbf{Stage 1:} HER2+ vs rest, fixed KW+RF $k=5$, threshold 0.5 (no inner CV tuning)
  \item Top features: chr17q12 ERBB2 amplicon ($H = 73.9$, $p = 9 \times 10^{-17}$) + adjacent chr17 loci ($H = 22$--$26$)
  \item HER2+ shows CN $= 2$ at these loci; non-HER2+ shows CN $\leq 0$ -- trivially separable
  \item BA = 1.000 in all 1000 outer folds; threshold 0.5 is in wide margin between class scores
  \item \textbf{Stage 2:} HR+ vs TN, one of the 4 pipeline variants with full inner CV tuning
  \item Feature selection and tuning on non-HER2+ samples only per outer fold
  \item Accesses chr12q/5q signal that 3-class ranking obscures
\end{itemize}

\subsection{Cross-validation scheme}
\label{sec:cv}

\begin{itemize}[leftmargin=*, nosep]
  \item Nested stratified repeated 5-fold CV \citep{wessels2005protocol, ambroise2002selection}
  \item Outer loop: unbiased BA estimation, preserves 32/36/32 proportions
  \item Flat: 50 repeats $\times$ 5 folds = 250 evaluations (seeds 1--50)
  \item Hierarchical: 200 repeats $\times$ 5 folds = 1000 evaluations (seeds 1001--1200)
  \item Inner loop: 5-fold CV on outer training set for hyperparameter tuning
  \item Feature selection INSIDE each inner fold (prevents leakage)
  \item Metric: balanced accuracy (macro-averaged recall) \citep{brodersen2010balanced}
  \item Best inner config selected $\to$ retrain on full outer train $\to$ evaluate on outer test fold
  \item See Figure~\ref{fig:workflow} for overview
\end{itemize}

\subsection{Plateau ensembling}
\label{sec:plateau-method}

\begin{itemize}[leftmargin=*, nosep]
  \item Problem: inner CV on $n \approx 11$ cannot reliably rank 800 configs; single-best is noise-dominated
  \item Solution: deploy all near-optimal configs instead of single best
  \item Pool inner CV scores across all outer folds and repeats
  \item Plateau threshold: best pooled mean $-$ 1 pooled std; retain up to 15 configs
  \item Prediction: average class probabilities across plateau ensemble, pick argmax
  \item Threshold and cap fixed before observing outer-fold test scores
  \item Leave-one-repeat-out: removing any repeat changes plateau by $<1$ of 15 configs
  \item Caveat: plateau gain may be modestly inflated by indirect inner/outer dependence; not quantified precisely
\end{itemize}

\subsection{Statistical comparison}
\label{sec:stats}

\begin{itemize}[leftmargin=*, nosep]
  \item Friedman test for overall pipeline differences
  \item Nadeau-Bengio corrected $t$-test \citep{nadeaubengio2003} for formal pairwise significance (accounts for fold overlap)
  \item Bootstrap CIs (10,000 resamples) on per-repeat mean differences for pragmatic distinguishability
  \item \cite{bouckaert2004evaluating}: standard paired tests overstate significance for repeated CV
  \item Do not use Wilcoxon (treats folds as independent, anti-conservative)
\end{itemize}

\subsection{Final model training}
\label{sec:final-model}

\begin{itemize}[leftmargin=*, nosep]
  \item Best pipeline: hierarchical EN+NMC with plateau ensembling
  \item Retrained on all 100 labelled samples
  \item Plateau membership from nested CV (no re-tuning)
  \item Merging map applied to validation data for consistent segment boundaries
  \item Accuracy estimate: $\sum_c (\hat{r}_c \times n_c / N) \times 57$ using training proportions as proxy
  \item Sensitivity: estimate ranges 46--48 under alternative plausible proportions; submit 48
\end{itemize}

% =============================================================================
\section{Results}
\label{sec:results}

\subsection{Flat 2x2 experiment}
\label{sec:results-flat}

\begin{itemize}[leftmargin=*, nosep]
  \item KW+RF best (BA = 0.791), KW+NMC worst (BA = 0.631)
  \item RF $\gg$ NMC regardless of feature selector (+16 pp for KW, +2.4 pp for EN)
  \item Feature selector matters less: KW vs EN only $\sim$2 pp difference for RF
  \item Nadeau-Bengio: only KW+NMC vs KW+RF significant after correction
  \item HER2+ recall = 1.00 in all 4 pipelines
  \item Bottleneck: TN recall = 0.62--0.64 (HR+/TN confusion)
  \item Motivates hierarchical decomposition (next subsection)
\end{itemize}

\subsection{Hierarchical classifier}
\label{sec:results-hierarchical}

\begin{itemize}[leftmargin=*, nosep]
  \item NMC $>$ RF reversal in Stage 2 (bootstrap CIs exclude zero)
  \item EN+NMC(plateau) best: BA2 = 0.759, combined BA = 0.840
  \item Nadeau-Bengio: all pairwise comparisons n.s.\ (conservative)
  \item Bootstrap CIs: 38/45 pairs distinguishable at 95\%
  \item TABLE 1: flat + hierarchical results combined (BA, BA2, per-class recalls, features)
  \item FIGURE 2: violin/box plot, flat vs hierarchical BA, significance brackets
\end{itemize}

\subsection{Plateau ensembling}
\label{sec:plateau}

\begin{itemize}[leftmargin=*, nosep]
  \item EN+NMC base BA2 = 0.729 $\to$ plateau BA2 = 0.759 (+3.0 pp)
  \item Gain scales with grid size: 800 configs $\to$ +3.0 pp; 32 configs $\to$ +1.4 pp
  \item Mechanism: averaging stabilises noisy single-best selection
  \item Leakage caveat: may be modestly inflated; leave-one-repeat-out suggests small effect
\end{itemize}

\subsection{Single best biomarker region}
\label{sec:biomarker}

\begin{itemize}[leftmargin=*, nosep]
  \item chr17q12 (ERBB2/TCAP locus), KW $H = 73.9$, $p = 9 \times 10^{-17}$
  \item Perfectly separates HER2+ from rest (Stage 1 uses 5 features from this locus)
  \item Cannot distinguish HR+ vs TN (different signal: chr12q, chr5q)
  \item Biology: ERBB2 amplification defines HER2+ subtype, targeted by trastuzumab \citep{slamon1987her2}
\end{itemize}

\subsection{Accuracy estimate}
\label{sec:estimate}

\begin{itemize}[leftmargin=*, nosep]
  \item Formula: $\sum_c (\hat{r}_c \times n_c / N) \times 57$ (not BA $\times$ 57)
  \item Using training proportions: $(0.32 \times 1.00 + 0.36 \times 0.78 + 0.32 \times 0.74) \times 57 = 47.8 \to 48$
  \item Sensitivity: uniform $\to$ 48, skewed $\to$ 46; submit 48
\end{itemize}

\subsection{Improvements after draft}
\label{sec:improvements}

\begin{itemize}[leftmargin=*, nosep]
  \item (a) Hierarchical architecture replacing flat 3-class
  \item (b) Plateau ensembling for deployment stability
  \item (c) Updated accuracy estimate with per-class formula
  \item (d) Quantify improvement: draft BA $\to$ final BA, delta in expected correct
  \item [TO BE FINALISED after draft feedback]
\end{itemize}

% =============================================================================
\section{Discussion}
\label{sec:discussion}

% ---------------------------------------------------------------------------
% CONDITIONALITY CLAIM (key framing for this section):
%
%   "Wessels' principle is conditional on multi-class structure. On the full
%    3-class aCGH problem, RF outperforms NMC. On the decomposed binary HR+/TN
%    subproblem, NMC outperforms RF. This was not tested in the original study."
%
% Do NOT write "Wessels was right" or "our results validate Wessels."
% ---------------------------------------------------------------------------

\begin{itemize}[leftmargin=*, nosep]
  \item Wessels' principle is CONDITIONAL on problem structure
  \item 3-class: RF $>$ NMC (contradicts Wessels on expression data)
  \item Binary HR+/TN: NMC $>$ RF (consistent with Wessels)
  \item HER2+ acts as confounder: trivial separability inflates RF advantage in 3-class
  \item Feature selector matters less than classifier in both settings (consistent with Wessels)
  \item Limitations: $n=100$, only 68 for Stage 2, discrete CN loses fine-grained signal
  \item $\sim$10 consistently misclassified samples, no KW features distinguish them -- suggests BA2 $\approx 0.76$ ceiling is irreducible noise in discrete representation
  \item Best biomarker: chr17q12 ERBB2 for HER2+; chr12q/5q for HR+/TN
  \item Future: multi-omic integration, larger cohorts, investigate misclassified samples
\end{itemize}

% =============================================================================
% ---- Author contributions (CRediT) ----
\section*{Author contributions}

% TODO: Fill in using CRediT taxonomy roles.

% ---- References ----
\bibliography{references}

\end{document}
